% Un arbre LSM (organisation de RocksDB, compaction par niveaux) : chemin d'une
% écriture et chemin d'une lecture.
\documentclass[border=6pt]{standalone}
\usepackage{sd}
\begin{document}
\begin{tikzpicture}[x=1mm, y=1mm, note/.append style={font=\sffamily\small},
    b/.style={soft=#1, minimum height=8mm, font=\sffamily\small, inner sep=2pt},
    sst/.style={draw=#1, fill=#1Bg, line width=0.7pt, minimum width=#1W, minimum height=6mm, inner sep=1pt, font=\sffamily\footnotesize},
    f/.style={draw=Damber, fill=DamberBg, line width=0.6pt, minimum height=6mm, inner sep=1pt, font=\sffamily\footnotesize}]
  \node[b=Dslate, minimum width=24mm] (w) at (0,30) {écriture};
  \node[b=Dorange, minimum width=26mm] (wal) at (0,12) {journal (WAL)};
  \node[b=Dblue, minimum width=34mm] (mt) at (40,30) {memtable\\{\footnotesize triée, en mémoire}};
  \node[b=Dblue, minimum width=34mm] (imm) at (84,30) {memtable figée};
  \draw[arr=Dslate] (w) -- node[num=Dslate, above=1mm] {2} (mt);
  \draw[arr=Dslate] (w) -- node[num=Dslate, left=1mm] {1} (wal);
  \draw[arr=Dblue] (mt) -- node[lbl, above] {pleine} (imm);
  % niveaux
  \node[leg, anchor=east] at (28,12) {L0};
  \foreach \x in {0,1,2} \node[f, minimum width=16mm] at (40+18*\x,12) {SSTable};
  \draw[arr=Dblue] (imm.south) -- node[lbl, right] {vidage (flush)} (84,15.5);
  \node[note, anchor=west] at (96,12) {fichiers qui se chevauchent};
  \node[leg, anchor=east] at (28,0) {L1};
  \foreach \x in {0,...,3} \node[f, minimum width=12mm] at (37+13*\x,0) {};
  \node[note, anchor=west] at (96,0) {une suite triée, taille cible T};
  \node[leg, anchor=east] at (28,-12) {L2};
  \foreach \x in {0,...,9} \node[f, minimum width=5.6mm] at (34+6.2*\x,-12) {};
  \node[note, anchor=west] at (96,-12) {une suite triée, 10 × T};
  \node[leg, anchor=east] at (28,-24) {L3};
  \node[f, minimum width=62mm] at (61,-24) {…};
  \node[note, anchor=west] at (96,-24) {100 × T};
  \draw[arr=Dgreen] (14,6) -- node[lbl, left, text=Dgreen] {compaction} (14,-28);
  \node[note, anchor=north, align=center] at (60,-32) {lecture d'une clé : memtable, puis chaque fichier de L0, puis un seul fichier par niveau ;\\le filtre de Bloom de chaque fichier écarte presque tous ceux qui n'ont pas la clé};
\end{tikzpicture}
\end{document}
